\begin{table}[t]
    \centering
    \caption{Comparison with existing code generation benchmarks. The \textbf{Task} column classifies benchmarks based on the granularity of the generated output. For instance, \textbf{Function-level Code Generation} indicates that the target output is confined to the scope of a single function, even if the input context originates from different files in a repository. FeatBench distinguishes itself by targeting \textbf{Feature Implementation} based on \textbf{Realistic Natural Language Requirements}, strictly devoid of code hints.}
    \resizebox{0.95\textwidth}{!}{
    \setlength{\tabcolsep}{10pt}
    \begin{tabular}{l|cccc}
        \toprule 
        \textbf{Benchmark} & \textbf{Date} & \textbf{Task} & \textbf{Data Curation}& \textbf{Code Hints}\\
        \midrule 
        HumanEval\cite{humaneval} & Jul, 2021 & Function-level Code Generation& Manual & Signatures  \\
        ClassEval\cite{classeval} & Aug, 2023 & Class-level Code Generation& Manual & Class Skeleton  \\
        CoderEval\cite{yu2024codereval} & Feb, 2024 & Function-level Code Generation& Manual & Signatures \\
        DevEval\cite{li2024deveval} & Aug, 2024 & Function-level Code Generation& Manual & Signatures  \\
        EvoCodeBench\cite{li2024evocodebench}& Dec, 2024 & Function-level Code Generation& Automatic & Signatures  \\
        FEA-Bench\cite{li2025fea} & Mar, 2025 & Feature-level Code Generation& Automatic& Signatures \\
        NoCode-bench\cite{deng2025nocode} & Jul, 2025 & Feature-level Code Generation& Automatic& Identifier Hints \\
        \midrule 
        \textbf{FeatBench} (Ours) & Feb, 2026 & \textbf{Feature-level Code Generation}& \textbf{Automatic} & \textbf{None}  \\
        \bottomrule
    \end{tabular}
    \label{tab:comparison} }
\end{table}